\section{Localisation $\Reop(s) = 1/2$: Four Equivalent Mechanisms}
\label{sec:localisation_critique}

The central object of this section is the PT-canonical explanation of the
critical line. Four apparently distinct mechanisms appear at different
stages of the programme. They are in fact four faces of a single
structure: the invariant measure of the persistence dynamics.

\subsection{Symplectic Planck Cell: $u_\star\,p_\star = 2\pi$}
\label{sec:cellule_planck}

The dilation operator $\HPTBK$ generates the scale transformations on
the Mellin coordinate. Its canonical Liouville measure is
$du\,dp_u / (2\pi)$, a direct consequence of the canonical commutator
$[u, p_u] = i$.

\begin{proposition}[PT Planck cell]
\label{prop:cellule_planck}
The unit cell of the canonical phase space is
\begin{equation}
\boxed{\;u_\star\, p_\star \;=\; 2\pi\;}
\label{eq:cellule_planck}
\end{equation}
with the symmetric choice
$u_\star = p_\star = \sqrt{2\pi} \approx 2.5066$ (fixed point of the
involution $(u, p_u) \leftrightarrow (p_u, u)$). At fixed energy $\gamma$,
the hyperbola $\HPTBK = \gamma$ enters the admissible space only
on the segment
\[
u \,\in\, [u_\star,\, \gamma/p_\star] \;=\; [\sqrt{2\pi},\, \gamma/\sqrt{2\pi}],
\]
and the canonical dynamical cut-off equals
\begin{equation}
\umax(\gamma) \;=\; \dfrac{\gamma}{\sqrt{2\pi}}.
\label{eq:umax}
\end{equation}
\end{proposition}

\begin{proposition}[Microcanonical semi-classical density]
\label{prop:densite_NPT}
With the cut-off~\eqref{eq:umax} and the double barrier imposed by
the involution $(u, p_u) \leftrightarrow (p_u, u)$, the microcanonical
quantisation of $\HPTBK$ in symmetric Weyl order produces the integrated
density
\begin{equation}
\NPT(\gamma) \;=\; \dfrac{\gamma}{2\pi}\,\log\dfrac{\gamma}{2\pi\,e} \,+\, 1.
\label{eq:NPT}
\end{equation}
\end{proposition}
\label{sec:densite_semi_classique}

\begin{proof}[Sketch]
The microcanonical area between the curve $\HPTBK = \gamma$ and the
barrier $u = u_\star$ in the quadrant $u, p_u > 0$ equals
$A(\gamma) = \int_{u_\star}^{\umax(\gamma)} \gamma/u\,du
= \gamma\,\log\!\big(\umax(\gamma)/u_\star\big)
= \gamma\,\log\!\big(\gamma/2\pi\big)$.
The double barrier (the mirror sector $u, p_u < 0$ obtained by the
involution) doubles this area. The semi-classical Weyl quantisation in
symmetric order gives $\NPT(\gamma) = A(\gamma)/\pi + \text{const.}$ The
constant is fixed to $+1$ by the double self-adjoint extension
of $\HPTBK$ at the boundaries (see \S\ref{sec:bord_antiperiodique}).
Direct algebraic computation gives~\eqref{eq:NPT}.
\end{proof}

\paragraph{Comparison with Riemann--von Mangoldt.}
The classical Riemann--von Mangoldt density \cite{Edwards1974} is
\begin{equation}
N_{\mathrm{RvM}}(\gamma) \;=\; \dfrac{\gamma}{2\pi}\,\log\dfrac{\gamma}{2\pi\,e} \,+\, \dfrac{7}{8} \,+\, O(1/\gamma).
\label{eq:NRvM}
\end{equation}
The functional form of~\eqref{eq:NPT} and~\eqref{eq:NRvM} is
identical. The constant gap
\begin{equation}
\NPT(\gamma) - N_{\mathrm{RvM}}(\gamma) \;=\; \dfrac{1}{8}
\label{eq:residu_un_huitieme}
\end{equation}
is verified to $10^{-15}$ on four decades, $\gamma \in [14, 10^{3}]$.
The $2\pi$ appearing in the density $N(\gamma)$ is the canonical
Liouville measure on $(u, p_u)$, consequence of the commutator
$[u, p_u] = i$; the shift $\Reop(s) = 1/2$ is derived from the Jacobian
of the multiplicative Haar measure (\S\ref{sec:haar_jacobien}). The residue
$1/8$ admits the cohomological overdetermination of
\S\ref{sec:surdetermination_un_huitieme}.

\subsection{Multiplicative Haar Jacobian}
\label{sec:haar_jacobien}

For a stationary mode $f(\tau, u) = e^{i\omega\tau}\,\psi(u)$ of
the d'Alembert operator on the Mellin coordinate (the Fisher metric
of $\Sigma_{\mathrm{pers}}$ being Lorentzian for
$\mu > \mu_c \approx 6.97$, \cite{PT_GeoSpectral}), the Mellin shift
under the multiplicative Haar measure
$d\pi_{\log} = \sum_p \delta_{\log p}\,d\log p$ produces a factor
$1/\sqrt{p} = e^{-u/2}$ in front of the transform (with $u = \log p$ one
has $1/p = e^{-u}$ and $1/\sqrt{p} = e^{-u/2}$; it is the square root,
that of the Haar density, which appears). More precisely, the pole of
the Mellin transform moves from $s = \omega$ to
\begin{equation}
\boxed{\;s \;=\; \dfrac{1}{2} + i\omega.\;}
\label{eq:haar_shift}
\end{equation}
This is the shift of the multiplicative Haar measure. The real part
$1/2$ is not postulated: it comes from the natural Jacobian of the
invariant Haar measure on the primes.

\subsection{Unitary Transformation $e^{su/2}$}
\label{sec:transformation_unitaire}

The Ruelle--PT self-adjointness theorem (\cite{PT_Foundations})
establishes that the discrete Ruelle--PT operator $L_s$ is self-adjoint for
the Gibbs measure $\mu_s(p) = p^{-s}$ via the detailed balance associated
with the antidiagonal symmetry $T_3 = T_3^{\top}$. In Mellin coordinate,
the Gibbs measure becomes $\mu_s(p)\,dp = e^{-su}\,du$. The unitary
transformation
\begin{equation}
U : f(u) \,\mapsto\, e^{-su/2}\,f(u)
\label{eq:U_unitaire}
\end{equation}
brings the operator $L_s$ on $L^{2}(\Reals, e^{-su} du)$ to
$L^{2}(\Reals, du)$ by translating $\HPTBK$ by $is/2$. The shift appears
literally in the exponent: it is the operator identification
between discrete detailed balance ($\Ttwo$ of the sieve) and continuous
self-adjointness of $\HPTBK$.

\subsection{Antiperiodic Boundary Condition}
\label{sec:bord_antiperiodique}

On $L^{2}([\umin, \umax], du)$ with $0 < \umin < \umax < \infty$, the
defect indices of $\HPTBK$ equal $(n_+, n_-) = (1, 1)$. By von
Neumann, there exists a $U(1)$ family of self-adjoint extensions
parameterised by
\begin{equation}
\psi(\umax) \;=\; e^{i\theta}\,\psi(\umin), \qquad \theta \in [0, 2\pi).
\label{eq:BC_general}
\end{equation}

\begin{proposition}[PT-canonical selection: antiperiodicity $\theta = \pi$]
\label{prop:antiperiodicite}
The involution
\[
\iota : u \,\mapsto\, \umax + \umin - u
\]
is the continuous projection of the antidiagonal
$T_3 = \bigl(\begin{smallmatrix}0 & 1\\ 1 & 0\end{smallmatrix}\bigr)$ of
the sieve $\bmod\,3$. It is the unique continuous involution preserving the
boundaries as orbit and is an isometry of $L^{2}$. The boundary term
$\iota\,\HPTBK\,\iota^{-1} + \HPTBK$ vanishes \textbf{only} for
$\theta = \pi$ (antiperiodic), which forces the spectral symmetry
$\{\gamma_n, -\gamma_n\}$.
\end{proposition}

\begin{proof}[Sketch]
The family~\eqref{eq:BC_general} is parameterised by
$\theta \in [0, 2\pi)$. The involution $\iota$ maps a self-adjoint
extension with parameter $\theta$ to the one with parameter $-\theta$.
Invariance under $\iota$ therefore imposes $\theta \equiv -\theta \pmod{2\pi}$,
i.e.\ $\theta \in \{0, \pi\}$. The case $\theta = 0$ (periodicity) gives a
symmetric spectrum but does not preserve the antisymmetry of $\HPTBK$ at
the boundaries; only $\theta = \pi$ satisfies
$\iota\,\HPTBK\,\iota^{-1} + \HPTBK = 0$ and produces the required
spectral symmetry $\{\gamma_n, -\gamma_n\}$. Verified numerically at
$5 \cdot 10^{-12}$ (machine precision).
\end{proof}

The continuous antiperiodicity is therefore the continuous image of
the discrete antidiagonality $T_3$ of the sieve, which is itself
the simplest expression of a non-trivial $2 \times 2$ stochastic matrix.
Discrete detailed balance and continuous self-adjointness translate
into each other through the unitary transformation~\eqref{eq:U_unitaire}.

\subsection{Unification}
\label{sec:localisation_unification}

The four mechanisms above manifest four facets of the same
persistence dynamics, with its multiplicative Haar invariant measure
and its $\Ttwo$ reversibility (double stochasticity). The links
M1$\leftrightarrow$M2 (Liouville measure $\to$ multiplicative Haar) and
M3$\leftrightarrow$M4 (Gibbs detailed balance $\to$ antiperiodicity
$T_3$) are explicitly constructed; the full link
M1$\leftrightarrow$M2$\leftrightarrow$M3$\leftrightarrow$M4 in the sense
of a strict mathematical equivalence requires a categorical
formulation not included here. The three roles of $\sper = 1/2$ listed
in \S\ref{sec:parametre} (arithmetic $\Tone$, geometric $\Tfive$,
spectral Haar) are operator-theoretically identified by this section.
This is the main structural contribution of the programme.

\begin{table}[ht]
\centering
\begin{tabular}{@{}lll@{}}
\toprule
\textbf{Mechanism} & \textbf{Description} & \textbf{PT origin}\\
\midrule
M1 --- Symplectic cell & $u_\star\,p_\star = 2\pi$ & $[u, p_u] = i$ (Liouville)\\
M2 --- Haar Jacobian & $du = dp/p$, shift $+1/2$ & $\Ttwo$, double stochasticity\\
M3 --- Unitary transformation & $U = e^{-su/2}$ & Gibbs detailed balance $\to$ SA\\
M4 --- Boundary condition & $\theta = \pi$ antiperiodic & $T_3$, antidiagonal of the sieve\\
\bottomrule
\end{tabular}
\caption{The four equivalent mechanisms for the critical line
$\Reop(s) = 1/2$.}
\label{tab:quatre_mecanismes}
\end{table}

